\documentclass{article}
\usepackage{graphicx}
\usepackage{amsmath}
\usepackage{amsfonts}
\usepackage{bookmark}
\usepackage{booktabs}
\usepackage{exercises}
\graphicspath{ {} }
\usepackage{indentfirst}
\usepackage{amsthm}
\usepackage{yhmath}
\usepackage{geometry}
\usepackage{placeins}
\usepackage{csquotes}
\usepackage{pgfplots}
\newcommand*{\oldneg}{\mathord{\sim}}
\newcommand\showdiv[1]{\overline{\smash{)}#1}}
\newcommand\Mydiv[1]{%
\smash{\raise.3ex\hbox{$\big)$}}\mkern-8mu
        \overline{\enspace\strut#1}}
\setlength{\tabcolsep}{5pt}
\theoremstyle{definition}
\newtheorem{definition}{Definition}

\setlength{\parindent}{.5in}
\setlength{\parskip}{12pt}
\renewcommand{\arraystretch}{2}
\geometry{margin=1in}
\newcommand{\?}{\stackrel{?}{=}}
\pgfplotsset{width=7cm,compat=1.18}\usepgfplotslibrary{statistics}

\title{5.1 Answers}
\begin{document}
\maketitle

\section*{1}
If $z$ varies directly with the product of $x$ and $y$, then $z$ is said to vary jointly with $x$ and $y$.

\section*{2}
Consider a set of data pairs. The set shows inverse variation iff the product of each of the two elements of each of the pairs equals, or is approximately equal to, some constant. 
Then it could be said that that set of data pairs shows inverse variation to that constant.

\section*{8}
Neither.

\section*{9}
Direct.

\section*{10}
Inverse.

\section*{11}
C

\section*{18}
\[ -\frac{1}{6}=-\frac{a}{12} \] \[ 2=a \] meaning, when $x=3$, \[ y=\frac{2}{3} \]

\section*{19}
\[ -7=\frac{a}{\frac{5}{3}} \] \[ a=-\frac{35}{3} \] meaning, when $x=3$, \[ y=-\frac{\frac{35}{3}}{3}=-\frac{35}{9} \]

\section*{23}
The set of data pairs here shows inverse variation. This is because the product of the two points in each data pair is equal to 85. 

\section*{29}
\[ 75=a\times 5\times -3 \] \[ a=-5 \] meaning, when $x=-4$ and $y=5$, \[ z=-15\times -4\times 5=300 \]

\section*{30}
A

\section*{31}
\[ x=\frac{ay}{z} \]

\section*{32}
\[ y=axz^2 \]

\section*{33}
\[ w=\frac{axz}{y} \]

\section*{34}
The error here is that the given equation shows when $z$ varies directly with the square root of $w$ and inversely with $x$ and the cube of $y$, which was not the question. Instead, 
the answer is $z=\frac{axy^3}{\sqrt{w}}$.

\section*{35}
Let $f(x)$ and $g(x)$ be a direct variation function and an inverse variation function, repsectively, and $h(x)$ be their sum. Thus:
\[ h(x)=a_1x+\frac{a_2}{x} \] for $a_1$ and $a_2$ the constants of variation for $f(x)$ and $g(x)$, respectively. Thus we would have 
$5=2a_1+\frac{a_2}{2}$. One possibility is taking $a_1$ to be $2$ and $a_2$ to be $2$, meaning $f(x)=2x$ and $g(x)=\frac{2}{x}$.

\section*{36}
Suppose $x$ varies inversely with $y$ and $y$ varies inversely with $z$. Thus:
\[ x=\frac{a_1}{y} \] \[ y=\frac{a_2}{z} \] for some constants $a_1$ and $a_2$, meaning \[ x=\frac{a_1}{\dfrac{a_2}{z}}=\frac{a_1}{a_2}z \] This means $x$ varies directly with $z$, 
in particular when $a=\frac{a_1}{a_2}$.

\section*{40}
(a): We know that $f=\frac{a\sqrt{T}}{Ld}$. Thus, if $f=262$, $T=670$, $L=62$, and $d=0.1025$, \[ 262=\frac{a\sqrt{670}}{62\times 0.1025} \]
\[ a\approx 64.325 \] meaning $f=\frac{64.325\sqrt{T}}{Ld}$.


(b): We know that $f=\frac{64.325\sqrt{T}}{Ld}$, meaning $f=\frac{64.325\sqrt{1629}}{201.6\times 0.49}\approx 26.282$.

\section*{41}
(a): We konw that $F$ varies jointly with $m_1$ and $m_2$ and inversely with the square of $d$. Thus $F=\frac{am_1m_2}{d^2}$.


(b): The universal gravitational constant is defined to be the constant of variation in part (a), that is the constant $a$. Using the information in the box, we have:
\[ 3.53\times 10^{22}=\frac{a\times 5.98\times 10^{24}\times 1.99\times 10^{30}}{(1.50\times 10^{11})^2} \]
\[ 3.53\times 10^{22}=\frac{11.9002a\times 10^{54}}{2.25\times 10^{22}} \]
\[ 3.53\times 10^{22}=\frac{11.9002}{2.25}a\times 10^{32} \]
\[ a\approx 0.6674\times 10^{-10} \]


(c): Because the product of the two masses is in the numerator of the expression in part (a), the gravitational force increases as the masses of the two objects increases and the 
distance between the objects stays constant. Likewise, the gravitational force decreases as the masses of the two objects stay constant and the distance between the objects increases.

\section*{42}
(a): Because the load $P$ varies jointly with the beam's width $W$ and the square of its depth $D$, and inversely with its unsupported length $L$, we know that 
$P=\frac{aWD^2}{L}$ for some constant $a$. Thus if $P=\frac{aWD^2}{L}$, then $\frac{a2WD^2}{2L}=\frac{aWD^2}{L}=P$ when the width and the length of the beam are doubled, meaning $P$ 
does not change in this case.


(b): If the width and the depth of the beam are double, then $P=\frac{a2W(2D)^2}{L}=\frac{8aWD^2}{L}$, meaning $P$ increases by a factor of $8$.


(c): If all three dimensions are doubled, then $P=\frac{a2W(2D)^2}{2L}=\frac{4aWD^2}{L}$, meaning $P$ increases by a factor of $4$.


(d): As per part (c), we could double all three dimensions of the beam to increase its load by a factor of 4. Thus by part (a), we could also
yet again double the width and the length of the beam (in fact, we could then double each of them the same number of arbitrary times), 
as the load would stay the same (that is, stay at the original load increased by a factor of 4).

\end{document}
